% Exercises for Chapter 8 (antithetic variates, moment matching and control variates).
% Statements only: solutions are kept privately.
% Unless stated otherwise: the test case S = K = 100, T = 1, r = 1%, y = 2%, sigma = 20%,
% S_T = S exp(nu T + sigma sqrt(T) Z) with nu = r - y - sigma^2/2 and Z ~ N(0,1).

\begin{exercise}[Antithetic digitals]\label{ex:cv-digital}
\easy Let $D=\1_{\{S_T>K\}}$ be a digital call in the Black--Scholes model, $p=\Q(S_T>K)$,
and $\rho$ the correlation of $D$ computed from $Z$ and from $-Z$. You may use
\noteref{prop:cv-antithetic} and \noteref{prop:cv-disjoint}.
\begin{enumerate}[label=(\alph*)]
  \item For $K$ at or above the median $Se^{\nu T}$ of $S_T$, show that
    $\rho=-p/(1-p)$, and give the standard-error ratio of $n$ antithetic pairs to $n$ plain
    draws and the efficiency at twice the cost.
  \item For $K$ below the median, compute $\E[D(Z)D(-Z)]$ and show that
    $\rho=-(1-p)/p$. Conclude that $\rho=-\min(p,1-p)/\max(p,1-p)$ for every $K$.
  \item What happens at $K=Se^{\nu T}$? Explain without any formula why the antithetic
    estimator is then exact.
  \item Evaluate $p$, $\rho$ and the efficiency for $K=90$ and $K=110$ in the test case.
\end{enumerate}
\end{exercise}

\begin{exercise}[When pairing hurts]\label{ex:cv-straddle}
\medium This is a computational exercise. Consider the discounted straddle
$f(Z)=e^{-rT}|S_T-K|$.
\begin{enumerate}[label=(\alph*)]
  \item Let $K=Se^{\nu T}$, the median of $S_T$, and $s=\sigma\sqrt T$. Show that
    $e^{rT}f(Z)/(Ks)\to|Z|$ in $L^2$ as $s\to0$, and deduce that the antithetic correlation
    tends to $1$. (You may use that correlations are continuous under $L^2$ convergence when
    the limits have positive variance.)
  \item With $10^5$ pairs, estimate the antithetic correlation $\rho(K)$ for strikes
    $K\in[70,130]$ and find the range of strikes on which antithetic sampling is worse than
    plain sampling with the same number of payoff evaluations. Explain the result with the
    even--odd decomposition of \noteref{prop:cv-antithetic}.
  \item At $K=100$, use instead the control $X=e^{-rT}(S_T-K)^2$. Compute its mean in closed
    form, and compare the standard errors of: plain Monte Carlo with $2n$ draws; $n$
    antithetic pairs; the control variate with estimated coefficient on $2n$ draws; and the
    control variate applied to the $n$ antithetic pair averages.
\end{enumerate}
\end{exercise}

\begin{exercise}[Moment matching is a control variate]\label{ex:cv-skewness}
\medium Let $Z_1,\dots,Z_n$ be i.i.d.\ $\mathcal N(0,1)$ and $Z'_i=(Z_i-\bar Z_n)/\tilde s_n$ the
log-matched sample of \noteref{def:cv-mm}. Take $f(z)=z^3$, so that
$\hat\mu^{\mathrm{mm}}_n=\frac1n\sum_iZ'^3_i$ estimates $\E Z^3=0$. You may use the strong
law, the central limit theorem and Slutsky's lemma.
\begin{enumerate}[label=(\alph*)]
  \item Show that $\hat\mu^{\mathrm{mm}}_n$ is exactly unbiased.
  \item Show that $\sqrt n\,\hat\mu^{\mathrm{mm}}_n\Rightarrow\mathcal N(0,6)$.
    \emph{Hint:} expand $\frac1n\sum_i(Z_i-\bar Z_n)^3$.
  \item Compare with the plain estimator $\frac1n\sum_iZ_i^3$ and with the control-variate
    estimator using the control $Z$ with its optimal coefficient. Why does the second matched
    moment, that is the control $Z^2$, not help here?
\end{enumerate}
\end{exercise}

\begin{exercise}[How many batches?]\label{ex:cv-mm-coverage}
\medium This is a computational exercise. Estimate the test-case call by log moment
matching with $n=4096$ draws split into $m\in\{4,16,64,256,1024\}$ separately matched
batches, and form the nominal $95\%$ confidence interval from the batch means with the
Student quantile on $m-1$ degrees of freedom.
\begin{enumerate}[label=(\alph*)]
  \item For each $m$, repeat the experiment $2000$ times and report the bias, the standard
    deviation of the estimator, the average estimated standard error and the coverage.
  \item Explain the results with \noteref{prop:cv-moment-matching} and
    \noteref{rem:cv-mm-as-cv}. Predict, without simulating, the coverage for $m=128$.
  \item Give a rule for choosing $m$ as a function of $n$.
\end{enumerate}
\end{exercise}

\begin{exercise}[Comparing a call and a put]\label{ex:cv-crn-parity}
\medium Let $C=e^{-rT}(S_T-K)^+$ and $P=e^{-rT}(K-S_T)^+$, both computed from the same
$Z$, and suppose we estimate $\E C-\E P$ by simulation.
\begin{enumerate}[label=(\alph*)]
  \item Show that $\Cov(C,P)=-\E[C]\,\E[P]$, and deduce that common random numbers
    \emph{increase} the variance of the estimated difference compared with independent draws.
  \item Using $S_T(-z)=(Se^{\nu T})^2/S_T(z)$, write $P(-Z)$ as a payoff in $S_T$ and explain why
    it is strongly correlated with $C$. Evaluate, for the test case, the variance per draw of
    the estimated difference with independent draws, with common random numbers, and with
    $C(Z_i)$ paired with $P(-Z_i)$ (use closed forms or quadrature).
  \item What is the best estimator of $\E C-\E P$? What does this suggest about using the
    put as a control variate for the call?
\end{enumerate}
\end{exercise}

\begin{exercise}[A fixed coefficient]\label{ex:cv-fixed-beta}
\easy Let $X,Y\in L^2$ with $\Var X>0$, and let $\beta^*=\Cov(X,Y)/\Var X$.
\begin{enumerate}[label=(\alph*)]
  \item Show that $\Var(Y-\beta X)<\Var Y$ if and only if $\beta$ lies strictly between $0$
    and $2\beta^*$.
  \item Take $\beta=1$ and $\Var X=\Var Y$. Show that the variance factor is $2(1-\rho)$, that
    the control helps only if $\rho>\frac12$, and that if the control doubles the cost per
    sample it improves efficiency only if $\rho>\frac34$.
  \item Let $\beta=\beta^*(1+\varepsilon)$, a coefficient with relative error $\varepsilon$. Show
    that the variance factor $\Var(Y-\beta X)/\Var Y$ is $1-\rho^2+\varepsilon^2\rho^2$. For
    $\rho=0.9$ and $\rho=0.999$, compare $\varepsilon=0$ with $\varepsilon=0.2$, and comment on how
    accurately the coefficient must be known when the correlation is high.
\end{enumerate}
\end{exercise}

\begin{exercise}[The price of estimating the coefficient]\label{ex:cv-lavenberg-welch}
\medium Let $(X_i,Y_i)$, $i\le n$, be i.i.d.\ bivariate normal with correlation $\rho$ and
$\Var Y=\sigma_Y^2$, let $\mu_X=\E X$ be known, and let $\hat\mu^{\mathrm{cv}}_n$ be the
control-variate estimator with the least-squares coefficient, as in
\noteref{prop:cv-estimated-beta}. You may use: $\bar X_n$ and
$\sum_i(X_i-\bar X_n)^2$ are independent, $\sum_i(X_i-\bar X_n)^2/\Var X\sim\chi^2_{n-1}$, and
$\E[1/\chi^2_k]=1/(k-2)$ for $k>2$.
\begin{enumerate}[label=(\alph*)]
  \item Show that, for $n\ge4$,
    \[
      \Var\hat\mu^{\mathrm{cv}}_n=\frac{\sigma_Y^2(1-\rho^2)}{n}\,\frac{n-2}{n-3}.
    \]
  \item For $\rho=0.5$, for which $n$ does the estimator beat the plain sample mean? And for
    $\rho=0.1$?
  \item Discuss what this suggests when the control-variate regression is run on many small
    batches.
\end{enumerate}
\end{exercise}

\begin{exercise}[The asymptotic square-root call]\label{ex:cv-asc}
\medium This is a computational exercise. The \emph{asymptotic square-root call} pays
$F(S_T)=2K\big[1-\sqrt{K/S_T}\,\big]^+$. You may use its Black--Scholes price
\[
  e^{-rT}2K\Big[N(d_-)-\sqrt{K/S}\,e^{-\nu T/2+\sigma^2T/8}\,N\big(d_--\tfrac12\sigma\sqrt T\big)\Big].
\]
\begin{enumerate}[label=(\alph*)]
  \item Show that $F$ vanishes for $S_T\le K$ and that its right derivative at the strike is
    $F'(K+)=1$, and explain why the at-the-money call $(S_T-K)^+$ is a natural control. What
    does the replication of $F$ by calls say about the coefficient?
  \item With $10^5$ draws, estimate the price by plain Monte Carlo, and by control variates
    with the discounted call as control, with $\beta=1$ and with the estimated coefficient.
    Test each estimate against the closed form.
  \item Compute the correlation and $\beta^*$ (by quadrature or with many draws), and the
    efficiencies of the two control-variate estimators at twice the cost. Why is the
    estimated coefficient so much better than the replication coefficient here?
\end{enumerate}
\end{exercise}
